\documentclass[11pt,a4paper]{article}

% ── Codificación, fuentes y lenguaje ───────────────────────────────────────────
\usepackage[utf8]{inputenc}
\usepackage[T1]{fontenc}
\usepackage[default,scale=0.95]{sourcesanspro}
\renewcommand{\familydefault}{\sfdefault}
\usepackage[spanish,es-noshorthands,es-tabla]{babel}

% ── Geometría y espaciado ─────────────────────────────────────────────────────
\usepackage[top=2.2cm, bottom=2.5cm, left=2.2cm, right=2.2cm, headheight=15pt]{geometry}
\usepackage{setspace}
\setstretch{1.18}
\usepackage{parskip}

% ── Matemáticas y unidades ────────────────────────────────────────────────────
\usepackage{amsmath}
\usepackage{amssymb}
\usepackage{siunitx}

% ── Circuitos ─────────────────────────────────────────────────────────────────
\usepackage[european, straightvoltages]{circuitikz}

% ── Tablas ────────────────────────────────────────────────────────────────────
\usepackage{booktabs}
\usepackage{tabularx}
\usepackage{array}
\usepackage{longtable}
\usepackage{colortbl}
\usepackage{enumitem}

% ── Colores, cajas e iconos ───────────────────────────────────────────────────
\usepackage[dvipsnames]{xcolor}
\usepackage{tcolorbox}
\tcbuselibrary{skins, breakable}
\usepackage{fontawesome5}
\usepackage{graphicx}

% ── Listados de código (dark-mode infográfico) ────────────────────────────────
%   Estilo base: fondo oscuro con números de línea. Los bloques lstlisting se
%   envuelven automáticamente en una caja tcolorbox redondeada.
\usepackage{listings}

% ── Secciones con barra de color (infografía) ────────────────────────────────
\usepackage{titlesec}
\titleformat{\section}
  {\normalfont\LARGE\bfseries\color{azulNoche}}
  {\colorbox{cyanNeon}{\color{fondoOscuro}\thesection}}{0.7em}{}
  [\vspace{2pt}\color{cyanNeon}\rule{\linewidth}{1.8pt}]
\titlespacing*{\section}{0pt}{24pt}{10pt}
\titleformat{\subsection}
  {\normalfont\large\bfseries\color{azulNoche}}
  {\textcolor{cyanNeon}{\thesubsection}}{0.7em}{}
  [\vspace{1pt}\color{grisLinea}\rule{\linewidth}{0.6pt}]
\titlespacing*{\subsection}{0pt}{14pt}{6pt}
\titleformat{\subsubsection}
  {\normalfont\normalsize\bfseries\color{azulNoche}}
  {\textcolor{cyanNeon}{\thesubsubsection}}{0.7em}{}
\titlespacing*{\subsubsection}{0pt}{10pt}{4pt}

% ── Cabeceras y pies de página ────────────────────────────────────────────────
\usepackage{fancyhdr}
\usepackage{lastpage}
\pagestyle{fancy}
\fancyhf{}
\renewcommand{\headrulewidth}{0pt}
\renewcommand{\footrulewidth}{0pt}
\fancyfoot[C]{%
  \begin{minipage}{\linewidth}
    \vspace{2pt}
    \color{cyanNeon}\rule{\linewidth}{1.2pt}\\[2pt]
    \centering\color{grisTexto}\footnotesize
    Página \thepage\ de \pageref{LastPage}
  \end{minipage}%
}

% ── Hiperenlaces ──────────────────────────────────────────────────────────────
\usepackage[colorlinks=true, linkcolor=azulNoche, urlcolor=cyanNeon]{hyperref}
\sloppy
\emergencystretch 3em

% ── Paleta de colores — tecnológico dark-mode ──────────────────────────────────
\definecolor{fondoOscuro}{HTML}{0D1B2A}     % Dark navy — fondo banda título
\definecolor{verdeTurquesa}{HTML}{004D40}    % Verde turquesa — bandas de marca
\definecolor{azulNoche}{HTML}{1B3A6B}       % Midnight blue — secciones, marcos
\definecolor{azulMedio}{HTML}{2A5298}       % Azul medio — variante de acento
\definecolor{cyanNeon}{HTML}{00C8E0}        % Cyan eléctrico — acentos primarios
\definecolor{verdeSignal}{HTML}{00B887}     % Verde señal — fortalezas, éxito
\definecolor{naranjaVivo}{HTML}{FF7043}     % Naranja vivo — mejoras, advertencias
\definecolor{rojoAlerta}{HTML}{E53935}      % Rojo alerta — errores
\definecolor{grisPapel}{HTML}{F4F6FA}       % Fondo tarjetas claras
\definecolor{grisLinea}{HTML}{C8D0E0}       % Bordes sutiles
\definecolor{grisTexto}{HTML}{6B7A99}       % Texto secundario
\definecolor{amarilloNota}{HTML}{FFB300}    % Notas / advertencias doradas

% Colores específicos para bloques de código (dark-mode)
\definecolor{codigoFondo}{HTML}{1E2D3D}      % Fondo del bloque de código
\definecolor{codigoComentario}{HTML}{7A8FA6} % Comentarios
\definecolor{codigoCadena}{HTML}{FF9D5C}     % Cadenas / strings
\definecolor{codigoKeyword}{HTML}{00C8E0}    % Palabras clave
\definecolor{codigoNumero}{HTML}{00E5A0}     % Números

% ── Banda de título: portada infográfica ──────────────────────────────────────
%   Diseño en 2 capas: banda de color (por defecto fondo oscuro) + franja de
%   acento cyan abajo. Uso: \bandaTitulo[color]{título}{subtítulo}.
%   El icono se puede cambiar con \renewcommand{\iconoBanda}{...} (FontAwesome).
\newcommand{\iconoBanda}{microchip}
\newcommand{\bandaTitulo}[3][fondoOscuro]{%
  \begin{tcolorbox}[
    enhanced,
    colback=#1, colframe=#1,
    boxrule=0pt, arc=8pt, outer arc=8pt,
    width=\linewidth,
    top=18pt, bottom=6pt, left=14pt, right=14pt,
    borderline south={3.5pt}{0pt}{cyanNeon},
  ]
    \begin{minipage}[c]{0.07\linewidth}
      \centering
      \textcolor{cyanNeon}{\fontsize{30}{30}\selectfont\faIcon{\iconoBanda}}
    \end{minipage}%
    \hspace{8pt}%
    \begin{minipage}[c]{0.88\linewidth}
      {\fontsize{20}{24}\selectfont\bfseries\color{white}#2}\par\vspace{5pt}%
      \textcolor{cyanNeon!80!white}{\small\faIcon{angle-right}~#3}%
    \end{minipage}
    \vspace{4pt}
  \end{tcolorbox}%
  \vspace{8pt}%
}

% ── Tarjetas de datos (infografía) ────────────────────────────────────────────
\tcbset{
  tarjetaDato/.style={
    enhanced, breakable,
    arc=6pt, outer arc=6pt,
    colback=grisPapel, colframe=azulNoche,
    boxrule=1pt,
    fonttitle=\bfseries\small, coltitle=white,
    attach boxed title to top left={yshift=-3mm, xshift=6mm},
    boxed title style={
      colback=azulNoche, colframe=azulNoche,
      arc=4pt, boxrule=0pt,
    },
    drop fuzzy shadow=azulNoche!30!white,
    top=8pt, bottom=8pt, left=10pt, right=10pt,
  },
  tarjetaIcono/.style={
    enhanced, breakable,
    arc=6pt, outer arc=6pt,
    colback=white, colframe=grisLinea, boxrule=0.8pt,
    fonttitle=\bfseries\small, coltitle=azulNoche,
    borderline west={3pt}{0pt}{cyanNeon},
    drop fuzzy shadow=azulNoche!25!white,
    top=6pt, bottom=6pt, left=10pt, right=10pt,
  },
  cajaTitulo/.style={
    enhanced, breakable,
    arc=6pt, outer arc=6pt,
    colback=azulNoche!10!grisPapel, colframe=azulNoche,
    boxrule=1pt,
    fonttitle=\bfseries, coltitle=white,
    attach boxed title to top left={yshift=-3mm, xshift=6mm},
    boxed title style={
      colback=azulNoche, colframe=azulNoche,
      arc=4pt, boxrule=0pt,
    },
    drop fuzzy shadow=azulNoche!25!white,
    top=8pt, bottom=8pt, left=10pt, right=10pt,
  },
  cajaContenido/.style={
    enhanced, breakable,
    arc=5pt, outer arc=5pt,
    colback=grisPapel, colframe=grisLinea,
    boxrule=0.8pt,
    fonttitle=\bfseries\small, coltitle=azulNoche,
    top=6pt, bottom=6pt, left=10pt, right=10pt,
  },
  cajaFortaleza/.style={
    enhanced, breakable,
    arc=6pt, outer arc=6pt,
    colback=verdeSignal!8!white, colframe=verdeSignal,
    boxrule=1pt,
    borderline west={4pt}{0pt}{verdeSignal},
    fonttitle=\bfseries\small, coltitle=verdeSignal!70!black,
    drop fuzzy shadow=verdeSignal!20!white,
    top=6pt, bottom=6pt, left=10pt, right=10pt,
  },
  cajaMejora/.style={
    enhanced, breakable,
    arc=6pt, outer arc=6pt,
    colback=naranjaVivo!8!white, colframe=naranjaVivo,
    boxrule=1pt,
    borderline west={4pt}{0pt}{naranjaVivo},
    fonttitle=\bfseries\small, coltitle=naranjaVivo!80!black,
    drop fuzzy shadow=naranjaVivo!20!white,
    top=6pt, bottom=6pt, left=10pt, right=10pt,
  },
  cajaRecomendacion/.style={
    enhanced, breakable,
    arc=6pt, outer arc=6pt,
    colback=cyanNeon!6!white, colframe=cyanNeon!60!azulNoche,
    boxrule=1pt,
    borderline west={4pt}{0pt}{cyanNeon},
    fonttitle=\bfseries\small, coltitle=azulNoche,
    drop fuzzy shadow=azulNoche!20!white,
    top=6pt, bottom=6pt, left=10pt, right=10pt,
  },
}

% ── Ícono + texto (encabezados de sección y elementos) ────────────────────────
\newcommand{\iconotexto}[2]{\textcolor{cyanNeon}{\faIcon{#1}}~#2}

% ── Listados de código: estilo dark + auto-envoltura ─────────────────────────
\lstdefinestyle{estiloCodigo}{
    backgroundcolor=\color{codigoFondo},
    basicstyle=\ttfamily\small\color{white},
    commentstyle=\color{codigoComentario}\itshape,
    stringstyle=\color{codigoCadena},
    keywordstyle=\color{codigoKeyword}\bfseries,
    numberstyle=\tiny\color{codigoComentario},
    numbers=left,
    numbersep=10pt,
    showstringspaces=false,
    breaklines=true,
    breakatwhitespace=false,
    tabsize=4,
    frame=none,
    captionpos=b,
    aboveskip=6pt,
    belowskip=4pt,
    xleftmargin=14pt,
    framexleftmargin=14pt,
    literate=
      {á}{{\'a}}1 {é}{{\'e}}1 {í}{{\'i}}1 {ó}{{\'o}}1 {ú}{{\'u}}1
      {Á}{{\'A}}1 {É}{{\'E}}1 {Í}{{\'I}}1 {Ó}{{\'O}}1 {Ú}{{\'U}}1
      {ñ}{{\~n}}1 {Ñ}{{\~N}}1 {ü}{{\"u}}1 {Ü}{{\"U}}1
      {¿}{{?`}}1 {¡}{{!`}}1
}
\lstdefinestyle{estiloPython}{
    style=estiloCodigo,
    language=Python,
}
\lstset{style=estiloCodigo}
\BeforeBeginEnvironment{lstlisting}{%
  \begin{tcolorbox}[
    enhanced, breakable,
    arc=6pt, outer arc=6pt,
    colback=codigoFondo, colframe=azulNoche,
    boxrule=1pt,
    drop fuzzy shadow=azulNoche!25!white,
    top=2pt, bottom=2pt, left=0pt, right=0pt,
  ]%
}
\AfterEndEnvironment{lstlisting}{\end{tcolorbox}}

% ── Cajas de contenido temáticas ──────────────────────────────────────────────
\newtcolorbox{cajaRecuerda}{
    enhanced, breakable,
    arc=6pt, outer arc=6pt,
    colback=azulNoche!10!grisPapel,
    colframe=azulNoche, boxrule=1pt,
    borderline west={4pt}{0pt}{cyanNeon},
    fonttitle=\bfseries\small\color{white},
    title={\faIcon{info-circle}~Recuerda},
    attach boxed title to top left={yshift=-3mm, xshift=6mm},
    boxed title style={colback=azulNoche, arc=4pt, boxrule=0pt},
    drop fuzzy shadow=azulNoche!25!white,
    left=10pt, right=10pt, top=8pt, bottom=8pt,
}

\newtcolorbox{cajaConcepto}{
    enhanced, breakable,
    arc=6pt, outer arc=6pt,
    colback=verdeSignal!8!white,
    colframe=verdeSignal, boxrule=1pt,
    borderline west={4pt}{0pt}{verdeSignal},
    fonttitle=\bfseries\small\color{white},
    title={\faIcon{lightbulb}~Concepto clave},
    attach boxed title to top left={yshift=-3mm, xshift=6mm},
    boxed title style={colback=verdeSignal!80!black, arc=4pt, boxrule=0pt},
    drop fuzzy shadow=verdeSignal!20!white,
    left=10pt, right=10pt, top=8pt, bottom=8pt,
}

\newtcolorbox{cajaEjemplo}{
    enhanced, breakable,
    arc=6pt, outer arc=6pt,
    colback=cyanNeon!5!white,
    colframe=cyanNeon!70!azulNoche, boxrule=1pt,
    borderline west={4pt}{0pt}{cyanNeon},
    fonttitle=\bfseries\small\color{fondoOscuro},
    title={\faIcon{code}~Ejemplo},
    attach boxed title to top left={yshift=-3mm, xshift=6mm},
    boxed title style={colback=cyanNeon, arc=4pt, boxrule=0pt},
    drop fuzzy shadow=azulNoche!20!white,
    left=10pt, right=10pt, top=8pt, bottom=8pt,
}

\newtcolorbox{cajaEjercicio}{
    enhanced, breakable,
    arc=6pt, outer arc=6pt,
    colback=azulMedio!7!white,
    colframe=azulMedio, boxrule=1pt,
    borderline west={4pt}{0pt}{azulMedio},
    fonttitle=\bfseries\small\color{white},
    title={\faIcon{pencil-alt}~Ejercicio},
    attach boxed title to top left={yshift=-3mm, xshift=6mm},
    boxed title style={colback=azulMedio, arc=4pt, boxrule=0pt},
    drop fuzzy shadow=azulNoche!20!white,
    left=10pt, right=10pt, top=8pt, bottom=8pt,
}

\newtcolorbox{cajaReto}{
    enhanced, breakable,
    arc=6pt, outer arc=6pt,
    colback=naranjaVivo!6!white,
    colframe=naranjaVivo, boxrule=1pt,
    borderline west={4pt}{0pt}{naranjaVivo},
    fonttitle=\bfseries\small\color{white},
    title={\faIcon{fire}~Reto},
    attach boxed title to top left={yshift=-3mm, xshift=6mm},
    boxed title style={colback=naranjaVivo, arc=4pt, boxrule=0pt},
    drop fuzzy shadow=naranjaVivo!20!white,
    left=10pt, right=10pt, top=8pt, bottom=8pt,
}

% ── Indicadores de dificultad ─────────────────────────────────────────────────
\newcommand{\facil}{\textcolor{verdeSignal}{\faIcon{circle}\,\textbf{Fácil}}}
\newcommand{\intermedio}{\textcolor{naranjaVivo}{\faIcon{circle}\,\textbf{Intermedio}}}
\newcommand{\dificil}{\textcolor{rojoAlerta}{\faIcon{circle}\,\textbf{Difícil}}}

% ─────────────────────────────────────────────────────────────────────────────

\renewcommand{\iconoBanda}{cat}

% ── Caja de consejos ─────────────────────────────────────────────────────────
\newtcolorbox{cajaTips}{
    enhanced, breakable,
    arc=6pt, outer arc=6pt,
    colback=amarilloNota!6!white,
    colframe=amarilloNota, boxrule=1pt,
    borderline west={4pt}{0pt}{amarilloNota},
    fonttitle=\bfseries\small\color{fondoOscuro},
    title={\faIcon{lightbulb}~Consejos para enseñar},
    attach boxed title to top left={yshift=-3mm, xshift=6mm},
    boxed title style={colback=amarilloNota, arc=4pt, boxrule=0pt},
    drop fuzzy shadow=amarilloNota!20!white,
    left=10pt, right=10pt, top=8pt, bottom=8pt,
}

% ── Banda de título personalizada para Scratch ───────────────────────────────
\newcommand{\bandaTituloScratch}{%
  \bandaTitulo[verdeTurquesa]{Scratch}{Programa. Juega. Aprende.}%
}

\begin{document}

% ══════════════════════════════════════════════════════════════════════════════
% PORTADA — Banda de título
% ══════════════════════════════════════════════════════════════════════════════
\bandaTituloScratch

\begin{center}
{\large\textcolor{grisTexto}{
  Guía completa para empezar a programar en \textbf{\textcolor{azulNoche}{Scratch}}
  con una adolescente de 14 años en \textbf{\textcolor{azulNoche}{GNU/Linux Mint}}.}}
\end{center}

% ══════════════════════════════════════════════════════════════════════════════
\section{\iconotexto{puzzle-piece}{¿Qué es Scratch?}}

\textbf{Scratch} es un lenguaje de programación \textbf{visual} desarrollado por el
\textbf{MIT Media Lab} (dirigido por Mitchel Resnick). En lugar de escribir código
textual, los programas se construyen encajando \textbf{bloques visuales} como piezas
de un rompecabezas.

\begin{cajaConcepto}
\textbf{Programar en Scratch es encajar bloques.} Cada bloque representa una
instrucción (mover, sonar, repetir, decidir). Te enseñan la \textbf{lógica de
programación} sin la barrera de la sintaxis, ideal para aprender pensamiento
computacional y creatividad digital.
\end{cajaConcepto}

Principales capacidades:
\begin{itemize}
  \item Crear \textbf{juegos}, animaciones e historias interactivas.
  \item Incorporar \textbf{sonidos}, \textbf{sprites} (personajes) y fondos personalizados.
  \item \textbf{Comunidad online} para compartir proyectos (como GitHub, pero para Scratch).
  \item Extensible con extensiones: \texttt{micro:bit}, Makey Makey, LEGO, etc.
\end{itemize}

Uso recomendado:
\begin{itemize}
  \item Enseñar \textbf{lógica y estructura de programación} sin sintaxis.
  \item Prototipos rápidos de interacción y animación.
  \item Educación STEM en escuelas y aprendizaje en casa.
\end{itemize}

% ══════════════════════════════════════════════════════════════════════════════
\section{\iconotexto{layer-group}{Versiones y plataformas}}

\begin{center}
\begin{tabularx}{\linewidth}{l l X}
\toprule
\textbf{Versión} & \textbf{Público} & \textbf{Plataformas} \\
\midrule
\textbf{ScratchJr}   & Niños 5--7 años        & Tablets (iOS, Android, Chromebook) \\
\textbf{Scratch 3.0} & 8--16+ años            & Navegador web, Windows, macOS, \textbf{Linux} \\
\bottomrule
\end{tabularx}
\end{center}

Nota: la \textbf{aplicación de escritorio} para GNU/Linux es \textbf{Scratch
Desktop}. También está disponible \textbf{Scratch 4.x} en versión de escritorio.

% ══════════════════════════════════════════════════════════════════════════════
\section{\iconotexto{wrench}{Instalación en GNU/Linux Mint}}

Opciones disponibles (ya instalada en nuestro equipo):

\begin{center}
\begin{tabularx}{\linewidth}{l X X}
\toprule
\textbf{Método} & \textbf{Ventaja} & \textbf{Desventaja} \\
\midrule
\textbf{AppImage} & Portátil, no modifica el sistema, funciona siempre &
Tan solo hay que descargarlo manualmente \\
\textbf{Flatpak}  & Se integra con el Gestor de Software de Mint &
Descarga algo pesada (runtime \SI{300}{\mega\byte}) \\
\textbf{Scratch online} & Sin instalar nada, listo para usar &
Requiere conexión y un navegador (Firefox/Chromium) \\
\bottomrule
\end{tabularx}
\end{center}

\begin{tcolorbox}[cajaRecomendacion, title={\faIcon{thumbs-up}~Recomendación: AppImage}]
Es el método más limpio para GNU/Linux Mint: no
depende de Snap (problemático en Mint) ni de Flatpak. Pasos rápidos:
\begin{enumerate}
  \item Descargar desde \url{https://scratch.mit.edu/download}.
  \item Darle permisos de ejecución con \texttt{chmod +x Scratch-*.AppImage}.
  \item Ejecutarlo con doble clic o desde terminal con `./Scratch-*.AppImage`.
\end{enumerate}
\end{tcolorbox}

% ══════════════════════════════════════════════════════════════════════════════
\section{\iconotexto{gamepad}{Plan de aprendizaje por fases}}

\subsection{Fase 1: Familiarización (semanas 1--2)}

\begin{center}
\renewcommand{\arraystretch}{1.3}
\begin{tabularx}{\linewidth}{l X c}
\toprule
\textbf{Proyecto} & \textbf{Qué aprende} & \textbf{Dificultad} \\
\midrule
Hola Mundo            & Bloques de movimiento, apariencia y sonido       & \facil \\
Animación básica      & Secuencias, bucles \texttt{repeat}, \texttt{esperar} & \facil \\
El gato bailarín      & Bucle \texttt{forever}, eventos al hacer clic en la bandera verde & \facil \\
\bottomrule
\end{tabularx}
\end{center}

\subsection{Fase 2: Interactividad (semanas 3--4)}

\begin{center}
\renewcommand{\arraystretch}{1.3}
\begin{tabularx}{\linewidth}{l X c}
\toprule
\textbf{Proyecto} & \textbf{Qué aprende} & \textbf{Dificultad} \\
\midrule
Pong clásico          & Teclado, colisiones, variables (puntos)          & \intermedio \\
Historia interactiva  & Condicionales \texttt{si/si-no}, diálogos         & \intermedio \\
Quiz / Trivia         & Lógica, comparaciones y retroalimentación         & \intermedio \\
\bottomrule
\end{tabularx}
\end{center}

\subsection{Fase 3: Juegos (semanas 5--8)}

\begin{center}
\renewcommand{\arraystretch}{1.3}
\begin{tabularx}{\linewidth}{l X c}
\toprule
\textbf{Proyecto} & \textbf{Qué aprende} & \textbf{Dificultad} \\
\midrule
Laberinto saltarín       & Detección de bordes, coordenadas X/Y             & \dificil \\
Plataformas              & Física básica, gravedad y clonación              & \dificil \\
Space Invaders           & Clonación masiva y dificultad progresiva         & \dificil \\
\bottomrule
\end{tabularx}
\end{center}

% ══════════════════════════════════════════════════════════════════════════════
\section{\iconotexto{book}{Recursos recomendados}}

\subsection{Scratch 3 Programming Playground — libro gratuito online}

\textbf{Al Sweigart} (No Starch Press). Gratis para leer en línea bajo licencia
Creative Commons, con proyectos completos y progresivos:
\url{https://inventwithscratch.com/book3}.

\begin{center}
\renewcommand{\arraystretch}{1.3}
\begin{tabularx}{\linewidth}{c X}
\toprule
\textbf{Capítulo} & \textbf{Proyecto / contenido} \\
\midrule
0 & Introducción y conceptos básicos \\
1 & Primeros pasos: interfaz, movimiento y eventos \\
2 & Líneas Arcoíris: bucles, colores y lápiz \\
3 & \textbf{Laberinto}: colisiones, coordenadas y condiciones \\
4 & \textbf{Baloncesto con gravedad}: física, variables y teclado \\
5 & \textbf{Brick Breaker pulido}: clonación, puntuación y niveles \\
6 & \textbf{Asteroides en el espacio}: enemigos múltiples y dificultad \\
7 & \textbf{Plataformas avanzado}: gravedad, saltos y \textit{scrolling} \\
\bottomrule
\end{tabularx}
\end{center}

\begin{itemize}
  \item Incluye \textbf{descargables} de sprites y fondos:
        \url{https://inventwithscratch.com/ScratchPlayground3_resources.zip}.
  \item Ejercicios al final de cada capítulo con retos de personalización.
  \item Proyectos en orden creciente de dificultad, ideales para 14 años.
\end{itemize}

\subsection{Curso gratuito en Udemy}

\textit{Scratch Game Programming}, también de \textbf{Al Sweigart}, sigue los
primeros 6 proyectos del libro: \url{https://www.udemy.com/scratch-game-programming/}.

\subsection{Manual para padres y profesores}

\textit{Scratch Class Handbook} (PDF gratuito, actualizado oct 2020):
\url{https://inventwithscratch.com/Scratch_Class_Handbook_v3.pdf}.
Contiene estrategias para conducir la sesión, tiempos y evaluación.

\subsection{Tutoriales oficiales de Scratch}

\begin{itemize}
  \item Tutoriales paso a paso: \url{https://scratch.mit.edu/tutorials}.
  \item Proyectos de la comunidad para inspirarse: \url{https://scratch.mit.edu/explore}.
  \item Información para padres: \url{https://scratch.mit.edu/parents}.
\end{itemize}

% ══════════════════════════════════════════════════════════════════════════════
\section{\iconotexto{calendar-alt}{Plan de estudio detallado (8 semanas)}}

\subsection{Semanas 1--2: Fundamentos}

\begin{center}
\renewcommand{\arraystretch}{1.3}
\begin{tabularx}{\linewidth}{c l X}
\toprule
\textbf{Día} & \textbf{Actividad (30 min)} & \textbf{Referencia} \\
\midrule
L & Interfaz de Scratch y crear cuenta        & Capítulo 1 \\
M & Primer sprite: movimiento con teclado     & Capítulo 1 \\
X & Añadir sonido y fondo                     & Capítulo 1 \\
J & Proyecto: ``Mi primer dibujo'' (lápiz)    & Capítulo 2 \\
V & Líneas arcoíris                           & Capítulo 2 \\
\bottomrule
\end{tabularx}
\end{center}

\subsection{Semana 3: Condicionales}

\begin{center}
\renewcommand{\arraystretch}{1.3}
\begin{tabularx}{\linewidth}{c l X}
\toprule
\textbf{Día} & \textbf{Actividad} & \textbf{Referencia} \\
\midrule
L & Colisiones básicas              & Capítulo 3 \\
M & Laberinto (parte 1: movimiento) & Capítulo 3 \\
X & Laberinto (parte 2: paredes)    & Capítulo 3 \\
J & Laberinto (parte 3: meta)       & Capítulo 3 \\
V & Personalizar el laberinto       & Capítulo 3 \\
\bottomrule
\end{tabularx}
\end{center}

\subsection{Semana 4: Variables y puntuación}

\begin{center}
\renewcommand{\arraystretch}{1.3}
\begin{tabularx}{\linewidth}{c l X}
\toprule
\textbf{Día} & \textbf{Actividad} & \textbf{Referencia} \\
\midrule
L & Concepto de variables              & Capítulo 4 \\
M & Baloncesto (parte 1: la pelota)    & Capítulo 4 \\
X & Baloncesto (parte 2: la canasta)   & Capítulo 4 \\
J & Baloncesto (parte 3: gravedad)     & Capítulo 4 \\
V & Baloncesto (parte 4: puntuación)   & Capítulo 4 \\
\bottomrule
\end{tabularx}
\end{center}

\subsection{Semana 5: Clonación}

\begin{center}
\renewcommand{\arraystretch}{1.3}
\begin{tabularx}{\linewidth}{c l X}
\toprule
\textbf{Día} & \textbf{Actividad} & \textbf{Referencia} \\
\midrule
L & Concepto de clones                  & Capítulo 5 \\
M & Brick Breaker (parte 1: la pelota)  & Capítulo 5 \\
X & Brick Breaker (parte 2: ladrillos)  & Capítulo 5 \\
J & Brick Breaker (parte 3: niveles)    & Capítulo 5 \\
V & Brick Breaker (parte 4: pulir)      & Capítulo 5 \\
\bottomrule
\end{tabularx}
\end{center}

\subsection{Semanas 6--7: Juegos avanzados}

\begin{center}
\renewcommand{\arraystretch}{1.3}
\begin{tabularx}{\linewidth}{l l X}
\toprule
\textbf{Semana} & \textbf{Proyecto} & \textbf{Referencia} \\
\midrule
6 & Asteroides en el espacio     & Capítulo 6 \\
7 & Plataformas saltarinas       & Capítulo 7 \\
\bottomrule
\end{tabularx}
\end{center}

\subsection{Semana 8: Proyecto libre}

\begin{center}
\renewcommand{\arraystretch}{1.3}
\begin{tabularx}{\linewidth}{X X}
\toprule
\textbf{Actividad} \\
\midrule
La alumna elige un juego o animación que le guste y lo completa \\
Se publica en la comunidad Scratch (\url{https://scratch.mit.edu}) \\
\bottomrule
\end{tabularx}
\end{center}

% ══════════════════════════════════════════════════════════════════════════════
\section{\iconotexto{youtube}{Canales de YouTube en español}}

\begin{center}
\renewcommand{\arraystretch}{1.3}
\begin{tabularx}{\linewidth}{l X}
\toprule
\textbf{Canal} & \textbf{Tipo de contenido} \\
\midrule
Scratch en Español            & Tutoriales paso a paso desde cero \\
Programación para Niños       & Juegos y animaciones guiados \\
CodeaKids                     & Proyectos creativos y lúdicos \\
Profesor Sheldon              & Cursos completos en español \\
\bottomrule
\end{tabularx}
\end{center}

% ══════════════════════════════════════════════════════════════════════════════
\section{\iconotexto{lightbulb}{Consejos para enseñar (14 años)}}

\begin{cajaTips}
\begin{enumerate}
  \item \textbf{Empieza con proyectos que ella elija}: ¿juegos, animaciones, música?
  \item No intentes ``enseñar programación'' directamente: deja que ella descubra
        cómo lograr lo que quiere.
  \item Sesiones de \textbf{30--45 minutos} máximo; la atención se pierde rápido.
  \item Usa proyectos de la comunidad como \textbf{punto de partida} (fork).
  \item \textbf{Celebra cada logro}: el primer sprite que se mueve es motivador.
  \item Alterna práctica y teoría corta: 10 minutos de concepto, 20 de construir.
\end{enumerate}
\end{cajaTips}

% ══════════════════════════════════════════════════════════════════════════════
\section{\iconotexto{clipboard-check}{Lista de verificación para empezar}}

\begin{center}
\renewcommand{\arraystretch}{1.3}
\begin{tabularx}{\linewidth}{l X}
\toprule
\textbf{Estado} & \textbf{Paso} \\
\midrule
\faCheckSquare & Scratch instalado en GNU/Linux Mint \\
$\square$      & Crear cuenta en \url{https://scratch.mit.edu} (guardar proyectos) \\
$\square$      & Abrir el libro gratuito: \url{https://inventwithscratch.com/book3} \\
$\square$      & Descargar recursos: \url{https://inventwithscratch.com/ScratchPlayground3_resources.zip} \\
$\square$      & Empezar el Capítulo 1 \\
\bottomrule
\end{tabularx}
\end{center}

% ══════════════════════════════════════════════════════════════════════════════
\section*{\iconotexto{link}{Referencias}}
\addcontentsline{toc}{section}{Referencias}

\begin{itemize}
  \item Scratch (sitio oficial): \url{https://scratch.mit.edu}
  \item Invent with Scratch (Al Sweigart): \url{https://inventwithscratch.com}
  \item Scratch 3 Programming Playground (libro online): \url{https://inventwithscratch.com/book3}
  \item Curso Audiovisual en Udemy: \url{https://www.udemy.com/scratch-game-programming/}
  \item Scratch Class Handbook (PDF): \url{https://inventwithscratch.com/Scratch_Class_Handbook_v3.pdf}
\end{itemize}

\end{document}
